\documentclass[11pt]{article}
\usepackage[margin=1in]{geometry}
\usepackage{amsmath,amssymb,amsthm}
\usepackage[T1]{fontenc}
\usepackage{lmodern}
\usepackage{graphicx}
\usepackage[hidelinks]{hyperref}
\usepackage{microtype}

\newtheorem{theorem}{Theorem}
\newtheorem{lemma}{Lemma}
\newtheorem{corollary}{Corollary}
\newcommand{\C}{\mathbb C}
\newcommand{\R}{\mathbb R}
\newcommand{\Dens}{\mathcal D}
\newcommand{\Tr}{\operatorname{Tr}}

\title{The sharp constant for rank-one coordinate maps and unconditional frame rescaling}
\author{Codex agent lane 11}
\date{July 21, 2026}

\begin{document}
\maketitle

\begin{center}
\fbox{\parbox{0.92\linewidth}{\textbf{Status.}
Full result likely valid, subject to human review.  The optimal complex
constant in Tselishchev's Theorems 2.2 and 2.3 is $4/\pi$; the optimal real
rescaling constant is $\pi/2$.}}
\end{center}

\section*{Source problem}

Tselishchev proves that there is a universal $K$ such that every bounded map
$\Phi:\ell_\infty^n\to B(H)$ with rank-one coordinate images satisfies
$\|\Phi\|_{\mathrm{cb}}\leq K\|\Phi\|$, and proves the equivalent finite
rescaling theorem for nonzero vector pairs \cite[Theorems 2.2--2.3]{Tsel}.
His proof gives $K=2$.  Remark 2.6(1), on PDF page 5, says that Gaussian
variables improve the constant and identifies evaluation of the best
possible $K$ as an interesting problem.

\begin{center}
\includegraphics[width=0.96\linewidth]{figures/open_problem_crop.png}
\end{center}

The frame-multiplier formulation originated in Balazs--Freeman--Popescu--
Speckbacher \cite[Conjecture 1.2]{BFPS}.  We solve the sharp-constant problem
in both of Tselishchev's finite formulations.

\section*{Definitions and main theorem}

For complex Hilbert spaces, let $K_{\mathrm{cb}}$ be the least $K$ such that
\[
 \|\Phi\|_{\mathrm{cb}}\leq K\|\Phi\|
\]
for all finite-dimensional $H$, all $n$, and all
$\Phi:\ell_\infty^n\to B(H)$ satisfying
$\operatorname{rank}\Phi(e_k)\leq1$.

Let $K_{\mathrm{res}}^{\C}$ be the least $K$ in the following nonzero-vector
rescaling statement.  Whenever
\begin{equation}\label{eq:unc}
 \left\|\sum_{k=1}^n \varepsilon_k m_k
       \langle u,y_k\rangle x_k\right\|
 \leq C\|u\|\qquad (u\in H,\ |\varepsilon_k|=1),
\end{equation}
with $m_k\neq0$ and $x_k,y_k\neq0$, there are factors
$c_k\overline{d_k}=m_k$ for which both $(c_kx_k)$ and $(d_ky_k)$ have
Bessel bound at most $KC$.  Define $K_{\mathrm{res}}^{\R}$ analogously over
real Hilbert spaces, using signs.

\begin{theorem}[Sharp constants]\label{thm:main}
The optimal constants are
\[
 K_{\mathrm{cb}}=K_{\mathrm{res}}^{\C}=\frac4\pi,
 \qquad
 K_{\mathrm{res}}^{\R}=\frac\pi2.
\]
In particular, the Gaussian improvement suggested in Remark 2.6(1) is
sharp in both the completely bounded-map and rescaling formulations.
\end{theorem}

\section*{Proof intuition}

The upper bound comes from the exact first absolute Gaussian moment.  The
operator norm of a rank-one coordinate map yields an $\ell_1$ estimate for
all scalar matrix coefficients.  Applying that estimate to two independent
circular Gaussian sums turns the $d$-level matrix coefficient into a sum of
rank-one trace norms and costs exactly
$(\mathbb E|\gamma|)^{-2}=4/\pi$.

For rescaling, the same Gaussian calculation bounds a Hellinger-type sum for
every pair of density matrices.  A convex minimization over logarithmic
weights balances the two frame operators coordinate by coordinate, making
both Bessel bounds equal to that controlled sum.

Sharpness comes from independent Haar unit vectors.  In dimension $d$, the
ordinary norm of the associated empirical rank-one map approaches the square
of the first absolute coordinate moment.  Regardless of the sample, a
$d$-level amplification is at least $1/d$.  The quotient tends to $4/\pi$.
The same samples obstruct better Bessel rescaling by taking traces of the two
frame operators.

\section*{The sharp upper bounds}

\begin{lemma}[Complex completely bounded upper bound]\label{lem:cbupper}
If $\Phi:\ell_\infty^n\to B(H)$ has rank at most one on every coordinate,
then
\[
 \|\Phi\|_{\mathrm{cb}}\leq\frac4\pi\|\Phi\|.
\]
\end{lemma}

\begin{proof}
Write
\[
 \Phi(a)u=\sum_{k=1}^n a(k)\langle u,y_k\rangle x_k,
 \qquad C=\|\Phi\|.
\]
Taking scalar matrix coefficients and aligning phases gives
\begin{equation}\label{eq:l1}
 \sum_{k=1}^n
 |\langle u,y_k\rangle|\,|\langle x_k,v\rangle|
 \leq C\|u\|\|v\|.
\end{equation}
Fix $m$ and vectors $u_1,\ldots,u_m,v_1,\ldots,v_m\in H$.  Substitute
$u=\sum_j\gamma_ju_j$ and $v=\sum_i\delta_iv_i$ in \eqref{eq:l1}, where
the two families are independent standard circular complex Gaussians,
normalized by $\mathbb E|\gamma|^2=1$.  Since
\[
 \mathbb E|\gamma|=\frac{\sqrt\pi}{2},
\]
independence and Cauchy--Schwarz give
\begin{align}\label{eq:tracebound}
 &\sum_{k=1}^n
 \left(\sum_j|\langle u_j,y_k\rangle|^2\right)^{1/2}
 \left(\sum_i|\langle x_k,v_i\rangle|^2\right)^{1/2} \\
 &\hspace{1cm}\leq \frac4\pi C
 \left(\sum_j\|u_j\|^2\right)^{1/2}
 \left(\sum_i\|v_i\|^2\right)^{1/2}.
\end{align}

An element of $M_m(\ell_\infty^n)$ is a tuple $(A_k)_{k=1}^n$ with norm
$\max_k\|A_k\|$.  In a matrix coefficient of the amplification, the term
belonging to $k$ is the trace pairing of $A_k$ with the rank-one matrix
\[
 B_k=\big(\langle u_j,y_k\rangle\langle x_k,v_i\rangle\big)_{i,j}.
\]
Its trace norm is precisely the $k$th product on the left side of
\eqref{eq:tracebound}.  Trace duality proves
\[
 \|\Phi_{(m)}((A_k))\|\leq\frac4\pi C\max_k\|A_k\|.
\]
Taking the supremum over $m$ proves the lemma.
\end{proof}

\begin{lemma}[Gaussian rescaling upper bound]\label{lem:resupper}
The complex and real nonzero-vector rescaling constants satisfy
\[
 K_{\mathrm{res}}^{\C}\leq\frac4\pi,
 \qquad
 K_{\mathrm{res}}^{\R}\leq\frac\pi2.
\]
\end{lemma}

\begin{proof}
Absorb the symbols and their phases into nonzero vectors $X_k,Y_k$ so that
\eqref{eq:unc} becomes
\begin{equation}\label{eq:reduced}
 \left\|\sum_k\varepsilon_k\langle u,Y_k\rangle X_k\right\|
 \leq C\|u\|.
\end{equation}
Put $P_k=X_k\otimes X_k$ and $Q_k=Y_k\otimes Y_k$.  For density matrices
$\rho,\sigma\in\Dens(H):=\{R\geq0:\Tr R=1\}$, apply the scalar consequence
of \eqref{eq:reduced} to two independent Gaussian vectors with covariances
$\sigma$ and $\rho$.  Exactly as above,
\begin{equation}\label{eq:hellinger}
 \sum_k\sqrt{\Tr(\rho P_k)\Tr(\sigma Q_k)}\leq L,
\end{equation}
where $L=4C/\pi$ in the complex case.  For real standard Gaussians,
$\mathbb E|g|=\sqrt{2/\pi}$, so $L=\pi C/2$.

For $s\in\R^n$, define the convex functions
\[
 A(s)=\lambda_{\max}\!\left(\sum_k e^{s_k}P_k\right),
 \qquad
 B(s)=\lambda_{\max}\!\left(\sum_k e^{-s_k}Q_k\right).
\]
Because all $P_k,Q_k$ are nonzero, $A+B$ is coercive and has a minimizer
$s^*$.  The variational formula for the top eigenvalue and the subgradient
optimality condition give maximizing density matrices $\rho,\sigma$ such
that, for every $k$,
\begin{equation}\label{eq:balance}
 e^{s_k^*}\Tr(\rho P_k)=e^{-s_k^*}\Tr(\sigma Q_k).
\end{equation}
Indeed, the subgradients of $A$ are convex combinations of the vectors
$(e^{s_k}\Tr(\rho P_k))_k$ from maximizing states; their convex combinations
are again represented by a maximizing density matrix, and similarly for
$B$.

Writing $a_k=\Tr(\rho P_k)$ and $b_k=\Tr(\sigma Q_k)$, equation
\eqref{eq:balance} and \eqref{eq:hellinger} yield
\[
 A(s^*)=\sum_k e^{s_k^*}a_k=\sum_k\sqrt{a_kb_k}\leq L,
 \qquad B(s^*)\leq L.
\]
Thus $\alpha_k=e^{s_k^*/2}$ makes $(\alpha_kX_k)$ and
$(\alpha_k^{-1}Y_k)$ both $L$-Bessel.  Restoring the absorbed symbols proves
the lemma.
\end{proof}

\section*{Finite Haar examples}

Let $S_{\mathbb F}^{d-1}$ be the unit sphere of $\mathbb F^d$, with Haar
probability measure, for $\mathbb F=\C$ or $\R$.  Define
\[
 h_{\mathbb F,d}=\mathbb E|\langle e_1,Z\rangle|,
 \qquad Z\sim\operatorname{Haar}(S_{\mathbb F}^{d-1}).
\]

\begin{lemma}[Haar moments]\label{lem:moments}
The exact moments are
\begin{align}
 h_{\C,d}&=\frac{\Gamma(d)\Gamma(3/2)}{\Gamma(d+1/2)},
 & d h_{\C,d}^2&\longrightarrow\frac\pi4,\label{eq:cmoment}\\
 h_{\R,d}&=\frac{\Gamma(d/2)}{\sqrt\pi\,\Gamma((d+1)/2)},
 & d h_{\R,d}^2&\longrightarrow\frac2\pi.\label{eq:rmoment}
\end{align}
\end{lemma}

\begin{proof}
In the complex case, $|\langle e_1,Z\rangle|^2$ has the
$\operatorname{Beta}(1,d-1)$ distribution.  Integrating its square root gives
the first formula.  In the real case, the corresponding distribution is
$\operatorname{Beta}(1/2,(d-1)/2)$, giving the second formula.  The limits
follow from $\Gamma(t+a)/\Gamma(t+b)\sim t^{a-b}$.
\end{proof}

\begin{lemma}[Finite empirical realization]\label{lem:empirical}
For every fixed $d$, every $\eta>0$, and either scalar field, there are
finitely many pairs of unit vectors $(x_k,y_k)_{k=1}^N$ such that
\begin{equation}\label{eq:uniformemp}
 \sup_{\|u\|=\|v\|=1}
 \frac1N\sum_{k=1}^N
 |\langle u,y_k\rangle|\,|\langle x_k,v\rangle|
 \leq h_{\mathbb F,d}^2+\eta.
\end{equation}
\end{lemma}

\begin{proof}
Take independent Haar pairs $(X_k,Y_k)$.  For fixed unit $u,v$, the expected
summand factors and equals $h_{\mathbb F,d}^2$.  The functions
\[
 (x,y)\longmapsto |\langle u,y\rangle|\,|\langle x,v\rangle|
\]
are uniformly bounded and are Lipschitz in $(u,v)$ with a common constant.
Choose a finite $\delta$-net of the compact product of unit spheres, apply the
law of large numbers simultaneously on the net, and then let $\delta$ be
small.  Some finite sample satisfies \eqref{eq:uniformemp}.
\end{proof}

For such a sample in the complex case, define
\begin{equation}\label{eq:samplemap}
 \Phi_N(a)u=\frac1N\sum_{k=1}^N
 a(k)\langle u,y_k\rangle x_k.
\end{equation}
Aligning phases shows that the left side of \eqref{eq:uniformemp} is exactly
$\|\Phi_N\|$.  Hence
\begin{equation}\label{eq:sampleordinary}
 \|\Phi_N\|\leq h_{\C,d}^2+\eta.
\end{equation}

\begin{lemma}[Amplification lower bound]\label{lem:amplower}
For every collection of complex unit-vector pairs in $\C^d$, the map
\eqref{eq:samplemap} satisfies
\[
 \|\Phi_N\|_{\mathrm{cb}}\geq\|\Phi_{N,(d)}\|\geq\frac1d.
\]
\end{lemma}

\begin{proof}
Let $(e_j)_{j=1}^d$ be the standard basis and put
$u_j=v_j=e_j/\sqrt d$.  These define unit vectors in the Hilbert direct sum
$(\C^d)^d$.  For each $k$, the rank-one coefficient matrix
\[
 B_k=\big(\langle u_j,y_k\rangle\langle x_k,v_i\rangle\big)_{i,j=1}^d
\]
has trace norm
\[
 \|B_k\|_1=
 \left(\sum_j|\langle u_j,y_k\rangle|^2\right)^{1/2}
 \left(\sum_i|\langle x_k,v_i\rangle|^2\right)^{1/2}
 =\frac1d.
\]
Choose $A_k\in M_d$ with $\|A_k\|\leq1$ whose trace pairing with $B_k$
(with the harmless transpose dictated by conventions) equals $\|B_k\|_1$.
Then $A=(A_k)\in M_d(\ell_\infty^N)$ has norm at most one, while the matrix
coefficient of $\Phi_{N,(d)}(A)$ between $(u_j)$ and $(v_i)$ equals
\[
 \frac1N\sum_{k=1}^N\|B_k\|_1=\frac1d.
\]
\end{proof}

\begin{lemma}[Rescaling lower bound]\label{lem:reslower}
Regard the sample as a multiplier with symbols $m_k=1/N$.  If
$c_k\overline{d_k}=1/N$ and the sequences $(c_kx_k)$ and $(d_ky_k)$ have
Bessel bounds $B_x,B_y$, then
\[
 \max\{B_x,B_y\}\geq\frac1d.
\]
The same statement holds over the real field.
\end{lemma}

\begin{proof}
Taking traces of the two frame operators gives
\[
 \sum_k|c_k|^2\leq dB_x,
 \qquad
 \sum_k|d_k|^2\leq dB_y,
\]
because all $x_k,y_k$ are unit vectors.  Therefore
\[
 1=\sum_k|c_kd_k|
 \leq\left(\sum_k|c_k|^2\right)^{1/2}
      \left(\sum_k|d_k|^2\right)^{1/2}
 \leq d\sqrt{B_xB_y}.
\]
\end{proof}

\section*{Sharpness and conclusion}

For the complex samples, Lemmas \ref{lem:empirical} and \ref{lem:amplower}
give
\[
 \frac{\|\Phi_N\|_{\mathrm{cb}}}{\|\Phi_N\|}
 \geq\frac{1}{d(h_{\C,d}^2+\eta)}.
\]
First let $\eta\downarrow0$ and then $d\to\infty$.  By
\eqref{eq:cmoment}, this proves $K_{\mathrm{cb}}\geq4/\pi$.  Together with
Lemma \ref{lem:cbupper},
\[
 K_{\mathrm{cb}}=\frac4\pi.
\]

For any rescaling of the same multiplier, Lemma \ref{lem:reslower} forces a
Bessel bound of at least $1/d$, while its unconditional constant is at most
$h_{\C,d}^2+\eta$.  Thus
$K_{\mathrm{res}}^{\C}\geq4/\pi$.  Lemma \ref{lem:resupper} gives equality.

Repeating the empirical construction on the real sphere, using
\eqref{eq:rmoment} and the real part of Lemma \ref{lem:reslower}, gives
$K_{\mathrm{res}}^{\R}\geq\pi/2$.  The real upper bound in Lemma
\ref{lem:resupper} again gives equality.  This completes the proof of
Theorem \ref{thm:main}.

\section*{Verification notes and scope}

\begin{itemize}
\item Every sharpness example is finite-dimensional and has nonzero vectors
and nonzero symbols, so it belongs to the exact nondegenerate formulation of
Tselishchev's Theorem 2.2.
\item The lower bound requires no numerical simulation.  Random sampling is
used only as a probabilistic existence proof for a finite family; uniformity
follows from compactness and a finite net.
\item The amplification witness is at level $d$ and is explicit once the
sample is fixed.  Each matrix $A_k$ is a polar trace-duality witness for a
rank-one coefficient matrix.
\item If one retains irrelevant indices with a zero vector and a nonzero
symbol, the exact $4/\pi$ bound need not be attained by finite factors; every
constant strictly larger than $4/\pi$ still works.  The source's quantitative
Theorem 2.2 explicitly uses the nonzero-vector formulation.
\item This packet determines the rank-one constant only.  The growth of the
rank-$r$ constants $K(r)$ from Remark 2.6(2) remains open.
\end{itemize}

\section*{Bounded novelty check}

The run registry, solution, attempt, and proof-gap indexes were searched for
arXiv:2508.02802, ``rank-one completely bounded'', ``best constant'', and
``$4/\pi$''.  The local 2000--present arXiv source corpus was searched for the
same combinations.  Exact web searches on July 21, 2026 included the source
title and author, both symbolic forms of $4/\pi$, rank-one coordinate maps,
Haar unit-vector examples, and forward references to arXiv:2508.02802.
Pisier's account of the little Grothendieck theorem \cite[Section 5]{Pisier}
was checked as the Gaussian background explicitly cited by the source.  No
source stating the sharp constant for Tselishchev's two formulations was
found; the source itself says the best constant is beyond its scope.

The mechanism is closely related to classical sharp Gaussian constants, so
the recommended human review should particularly check whether this precise
rank-one-map specialization has appeared under different operator-ideal or
tensor-norm terminology.  Mathematically, the main review points are the
uniform empirical realization and the trace-pairing amplification witness.

\begin{thebibliography}{9}
\bibitem{BFPS}
P. Balazs, D. Freeman, R. Popescu, and M. Speckbacher,
\emph{Quantitative bounds for unconditional pairs of frames},
arXiv:2212.00947 (2022).

\bibitem{Pisier}
G. Pisier,
\emph{Grothendieck's theorem, past and present},
Bull. Amer. Math. Soc. (N.S.) \textbf{49} (2012), 237--323;
arXiv:1101.4195.

\bibitem{Tsel}
A. Tselishchev,
\emph{Rescaling of unconditional Schauder frames in Hilbert spaces and
completely bounded maps}, arXiv:2508.02802 (2025).
\end{thebibliography}

\end{document}
